\subsection{XGBoost Regularization Sweep}

As noted in Section~5.2.3, an initial deeper, less-regularized XGBoost configuration exhibited severe overfitting on the training set. Table~\ref{tab:xgb_regularization} contrasts this initial configuration with the shallow-tree, L1/L2-regularized configuration selected via the manual regularization sweep.

\renewcommand{\arraystretch}{1.4}
\setlength{\tabcolsep}{8pt}
\begin{table}[H]
\centering
\caption{XGBoost Regularization Sweep (EER, \%)}
\label{tab:xgb_regularization}
\begin{tabular}{l c c c}
\hline
\textbf{Configuration} & \textbf{Training EER} & \textbf{Development EER} & \textbf{Train--Dev Gap} \\
\hline
Initial (deep, unregularized) & \textit{TBD} & \textit{TBD} & \textit{TBD} \\
Final (shallow, L1/L2-regularized, selected) & \textit{TBD} & \textit{TBD} & \textit{TBD} \\
\hline
\end{tabular}
\end{table}

The initial configuration is expected to show near-zero training EER alongside a substantially higher development EER, indicating memorization of the training set rather than generalizable spoof-detection behavior; the regularized configuration is expected to narrow this gap, at the cost of a small increase in training EER, and is used for the XGBoost results reported in Section~6.4.
